Considera les funcions $f(x)=x^2-1$ i $g(x)=x+1$.

\begin{apartats}

\apartat[1,25]{0,75}
Troba els punts de tall de les gràfiques de $f$ i $g$, i una primitiva de $g(x)-f(x)$.

\begin{solucio}
$x^2-1=x+1$, és a dir, $x^2-x-2=0$: $x=-1$ i $x=2$, als punts $(-1,0)$ i $(2,3)$.
$g(x)-f(x)=-x^2+x+2$, amb primitiva $-\dfrac{x^3}3+\dfrac{x^2}2+2x$.
\end{solucio}

\begin{tria}{area-corbes}
\itemtria{area}{1}{1,25}
Calcula l'àrea de la regió tancada per les gràfiques de $f$ i $g$.

\begin{solucio}
Entre $-1$ i $2$, $g\ge f$ (per exemple, a $x=0$, $g=1>f=-1$):
\[
\int_{-1}^2(-x^2+x+2)\,dx=\left[-\frac{x^3}3+\frac{x^2}2+2x\right]_{-1}^2
=\frac{10}3-\left(-\frac76\right)=\frac92
\]
unitats quadrades.
\end{solucio}

\itemtria{tres-talls}{1}{1,25}
Calcula l'àrea de la regió tancada per les gràfiques de $y=x^3$ i $y=x$.

\begin{solucio}
$x^3=x$ dona $x=-1$, $0$ i $1$: hi ha dues regions, i la funció que queda per sobre canvia. A $[-1,0]$,
$x^3\ge x$, i a $[0,1]$, $x\ge x^3$:
\[
\int_{-1}^0(x^3-x)\,dx+\int_0^1(x-x^3)\,dx=\frac14+\frac14=\frac12
\]
unitats quadrades.
\end{solucio}
\end{tria}

\begin{nomesllarg}
\apartat{0,75}
Calcula l'àrea de la regió tancada per la paràbola $y=x^2$ i la recta $y=4$.

\begin{solucio}
Es tallen a $x=\pm2$, i entre aquests valors la recta és per sobre:
$\displaystyle\int_{-2}^2(4-x^2)\,dx=\left[4x-\frac{x^3}3\right]_{-2}^2=\frac{16}3+\frac{16}3=\frac{32}3$
unitats quadrades.
\end{solucio}
\end{nomesllarg}

\end{apartats}
